\section{Open Questions y direcciones a largo plazo}
\subsection{Interpretabilidad y regulación}
Weston considera que ``the transformer itself does reasoning'', pero también nos advierte que ``interpreting the neural network isn't easy'' (SRT 3204). Por ello, una de las direcciones a largo plazo es construir un pipeline de razonamiento trazable, de manera que el verified response, además del final answer, también pueda producir un trace para uso regulatorio.

\begin{table}[H]
\centering
\begin{tabular}{p{5cm}p{10cm}}
\toprule
Problema & Descripción \\
\midrule
Brecha de interpretabilidad & ¿Cómo extraer un reasoning trace del verified response para que un humano lo revise? \\
Atribución de responsabilidad & Cuando el meta judge se equivoca, ¿cómo rastrear si la falla proviene del reward model, del judge o de una operación incorrecta del prompt? \\
Señales multiagente & En un escenario con múltiples agentes (bot ↔ judge ↔ user), ¿cómo garantizar que la reward signal no sea objeto de adversarial exploitation? \\
\bottomrule
\end{tabular}
\caption{Open Questions para uso posterior en alignment y auditoría.}
\end{table}

\subsection{Interacción con sistemas externos}
Los modelos que se automejoran en el futuro necesitan ``Reason by actually interacting with people in the world, internet or itself'' (SRT 3282-3284). Esto incluye:
\begin{itemize}
    \item Coordinarse con otros agentes (share verified responses, cross judge);
    \item Recabar additional context mediante API / internet, para después autoverificarse dentro de Self-Instruct;
    \item Buscar automáticamente la intervención de un agente externo al detectar una hallucination (self-aware behaviour).
\end{itemize}

\subsection{Resumen de la sección}
Open Questions ofrece una dirección para el alignment y la regulación; la interacción con el mundo exterior es el siguiente paso para construir un sistema self-improving más autónomo, más autoconsciente y también más seguro.

